
\section{Implementation}
\label{sec:implementation}


Put an introductory paragraph here that gives the reader an overview of what's coming. If there are multiple subsections, say in a few words or a sentence something about each subsection.

\subsection{Part 1}

State  the objective for the implementation (2-3 sentences).

Describe the implementation (2-3 sentences, or more if needed). Feel free to use coding examples as needed, such as that shown in Listing~\ref{listing:stencil-core}.

\begin{lstlisting}[caption={Stencil computation in 2D: performs sum of product of nearby pixels with weights.},label={listing:stencil-core}, name=stencil-core, float=h, style=mystyle,language=C++]
float smoothPixel(Si, Sj, S, R, weights) {
    // compute the weight sum of pixels nearby
    // this code doesn't handle edge conditions
    // and assumes sum of weights[i,j] = 1.0 
    float sum = 0.0;
    for (int j=0; j<R; j++)
        for (int i=0; i<R; i++)
           sum += weights[i,j]*S[Si+i,Sj+j]
    return sum; }
\end{lstlisting}

\subsection{Part 2}

Repeat the concepts/writing motif from above as needed for your implementations.

Here are some math examples if you need them for your work:



We can write the \emph{quantum state} encoded by a single qubit as
\begin{equation}
    \ket{\psi} \in \{ x : x \in \mathbb{C}^2, \norm{x} = 1 \}
\end{equation}
so that it is a complex-valued vector in 2 dimensions with unit norm. The notation $\ket{\psi}$ is known as \emph{Dirac notation} and is the standard notation for writing quantum states. (We will see later in this section how to visualize a 1-qubit quantum state using a Bloch sphere.) The quantum states
\begin{equation}
\ket{0} = \begin{pmatrix}
1 \\
0
\end{pmatrix}, 
\ket{1}= \begin{pmatrix}
0 \\
1
\end{pmatrix}
\end{equation}
correspond to the classical bit states of 0 and 1. As the astute reader may have observed, these two states form a small subset of possible quantum states. Nevertheless, we still have a tractable method to describe quantum states for the purposes of computation.

Classical sorting algorithms have $\mathcal{O}(N^2)$ complexity in the worst case.

As seen in Eq.~\ref{eq:vmm}, the vector-matrix multiplication algorithm computes $y = Ax$ where $x, y$ are column vectors of length $N$ and $A$ is a matrix of size $N \times N$:

\begin{equation}
y = \begin{bmatrix}
y_0 \\
y_1
\end{bmatrix}, 
x = \begin{bmatrix}
x_0 \\
x_1
\end{bmatrix},
A = \begin{bmatrix}
A_{0,0} & A_{0,1} \\
A_{1,0} & A_{1,1} 
\end{bmatrix}
\label{eq:vmm}
\end{equation}
